%% start of file `template.tex'.
%% Copyright 2006-2015 Xavier Danaux (xdanaux@gmail.com).
%
% This work may be distributed and/or modified under the
% conditions of the LaTeX Project Public License version 1.3c,
% available at http://www.latex-project.org/lppl/.
\documentclass[11pt,letter,sans,colorlinks,linkcolor=gray]{moderncv}        % possible options include font size ('10pt', '11pt' and '12pt'), paper size ('a4paper', 'letterpaper', 'a5paper', 'legalpaper', 'executivepaper' and 'landscape') and font family ('sans' and 'roman')
\pdfoutput=1
% moderncv themes
\moderncvstyle{classic}                             % style options are 'casual' (default), 'classic', 'banking', 'oldstyle' and 'fancy'
\moderncvcolor{blue}                               % color options 'black', 'blue' (default), 'burgundy', 'green', 'grey', 'orange', 'purple' and 'red'
%\renewcommand{\familydefault}{\sfdefault}         % to set the default font; use '\sfdefault' for the default sans serif font, '\rmdefault' for the default roman one, or any tex font name
%\nopagenumbers{}     
% uncomment to suppress automatic page numbering for CVs longer than one page
\usepackage{lastpage}
\rfoot{\addressfont\itshape\textcolor{gray}{\thepage/\pageref{LastPage}}}
% character encoding
\usepackage[utf8]{inputenc}                       % if you are not using xelatex ou lualatex, replace by the encoding you are using
%\usepackage{CJKutf8}                              % if you need to use CJK to typeset your resume in Chinese, Japanese or Korean

%%%%%%%%%%%%%%%

% PUBLICATIONS

% \usepackage[backend=biber,style=phys,biblabel=brackets,sorting=ndymdt,eprint=true]{biblatex}
\usepackage[backend=biber,style=phys,biblabel=brackets,sorting=ndymdt,eprint=true]{biblatex}
\addbibresource{data/publications.bib}

\DeclareFieldFormat[misc]{title}{\mkbibquote{#1}}

\defbibenvironment{publist}
  {\list
     {\hspace{-4em}\printtext[labelnumberwidth]{\printfield{labelnumber}}}
     {\setlength{\leftmargin}{21.25mm}%
      \setlength{\itemsep}{\bibitemsep}%
      \setlength{\parsep}{\bibparsep}
      }%
      \renewcommand*{\makelabel}[1]{\hss##1}
      }
  {\endlist}
  {\item}

% sort in reverse date order
% https://tex.stackexchange.com/questions/46868/biblatex-sorting-by-date
\DeclareSortingTemplate{ndymdt}{
  \sort{
    \field{presort}
  }
  \sort[final]{
    \field{sortkey}
  }
  % \sort{
  %   \field{sortname}
  %   \field{author}
  %   \field{editor}
  %   \field{translator}
  %   \field{sorttitle}
  %   \field{title}
  % }
  \sort[direction=descending]{
    \field{sortyear}
    \field{year}
    \literal{9999}
  }
  \sort[direction=descending]{
    \field[padside=left,padwidth=2,padchar=0]{month}
    \literal{99}
  }
  \sort[direction=descending]{
    \field[padside=left,padwidth=2,padchar=0]{day}
    \literal{99}
  }
  \sort{
    \field{sorttitle}
  }
  \sort[direction=descending]{
    \field[padside=left,padwidth=4,padchar=0]{volume}
    \literal{9999}
  }
}

% reverse numbering
% Count total number of entries in each refsection
\AtDataInput{%
  \csnumgdef{entrycount:\therefsection}{%
    \csuse{entrycount:\therefsection}+1}}

% Print the labelnumber as the total number of entries in the
% current refsection, minus the actual labelnumber, plus one
\DeclareFieldFormat{labelnumber}{\mkbibdesc{#1}}    
\newrobustcmd*{\mkbibdesc}[1]{%
  \number\numexpr\csuse{entrycount:\therefsection}+1-#1\relax}

%% replace my name with bold
% can do this, or can manually do it in bib file with
% @preamble{ "\newcommand{\jtiosue}{\mkbibbold{J. T. Iosue}}" }
\DeclareSourcemap{
  \maps[datatype=bibtex]{
    \map[overwrite]{
      \step[fieldsource=author, match={Iosue, Joseph T.}, replace={\\mkbibbold{J. T. Iosue}}]
      \step[fieldsource=author, match={Iosue, Joseph}, replace={\\mkbibbold{J. T. Iosue}}]
      \step[fieldsource=author, match={Joseph T. Iosue}, replace={\\mkbibbold{J. T. Iosue}}]
      \step[fieldsource=author, match={Joseph Iosue}, replace={\\mkbibbold{J. T. Iosue}}]
      \step[fieldsource=author, match={J. T. Iosue}, replace={\\mkbibbold{J. T. Iosue}}]
      \step[fieldsource=author, match={J. Iosue}, replace={\\mkbibbold{J. T. Iosue}}]
    }
  }
}
% biblatex wants integer months
% https://tex.stackexchange.com/questions/438396/regex-replacement-for-biblatex-biber-to-eliminate-the-warning-month-not-integer
% https://tex.stackexchange.com/questions/275564/biber-configuration-file-for-shortening-journal-names
% https://tex.stackexchange.com/questions/115283/how-to-use-biber-to-convert-integer-to-bibtex-month-macro
\DeclareSourcemap{
  \maps[datatype=bibtex]{
    \map[overwrite]{
      \step[fieldsource=month, match=\regexp{\A(j|J)an(uary)?\Z}, replace=1]
      \step[fieldsource=month, match=\regexp{\A(f|F)eb(ruary)?\Z}, replace=2]
      \step[fieldsource=month, match=\regexp{\A(m|M)ar(ch)?\Z}, replace=3]
      \step[fieldsource=month, match=\regexp{\A(a|A)pr(il)?\Z}, replace=4]
      \step[fieldsource=month, match=\regexp{\A(m|M)ay\Z}, replace=5]
      \step[fieldsource=month, match=\regexp{\A(j|J)un(e)?\Z}, replace=6]
      \step[fieldsource=month, match=\regexp{\A(j|J)ul(y)?\Z}, replace=7]
      \step[fieldsource=month, match=\regexp{\A(a|A)ug(ust)?\Z}, replace=8]
      \step[fieldsource=month, match=\regexp{\A(s|S)ep(tember)?\Z}, replace=9]
      \step[fieldsource=month, match=\regexp{\A(o|O)ct(ober)?\Z}, replace=10]
      \step[fieldsource=month, match=\regexp{\A(n|N)ov(ember)?\Z}, replace=11]
      \step[fieldsource=month, match=\regexp{\A(d|D)ec(ember)?\Z}, replace=12]
    }
  }
}


%%%%%%%%%%%%%%%

% adjust the page margins
\usepackage[right=1.9cm,left=1.7cm,bottom=2.5cm,top=2.6cm]{geometry}
\setlength{\hintscolumnwidth}{1.8cm}                % if you want to change the width of the column with the dates
\setlength{\makecvheadnamewidth}{8cm}           % for the 'classic' style, if you want to force the width allocated to your name and avoid line breaks. be careful though, the length is normally calculated to avoid any overlap with your personal info; use this at your own typographical risks...
% \usepackage{mathpazo}

% \definecolor{c1}{rgb}{0.858, 0.188, 0.478}
\definecolor{c1}{HTML}{B00000}
\usepackage{etaremune}
%To allow resumption after page break and such:
\usepackage{etoolbox} %% For \AtEndEnvironment

\makeatletter %% <- change @ so that it can be used in macro names
  %% Define the resume key for etaremune:
  \define@boolkey[EM]{etaremune}{resume}[true]{}
  \presetkeys[EM]{etaremune}{resume=false}{} %% <- false by default

  %% Increase starting value of previous etaremune environment if resuming
  \AtEndEnvironment{etaremune}{%
    \ifEM@etaremune@resume
      \EM@resumewrite{\csname EM@prevlist@\@roman{\@enumdepth}\endcsname}{\EM@currlist}%
    \fi
    \expandafter\xdef\csname EM@prevlist@\@roman{\@enumdepth}\endcsname{\EM@currlist}%
  }
  \newcommand*\EM@resumewrite[2]{% %% Expand arguments and then call EM@resume@write@
    \begingroup
      \edef\temp{\noexpand\EM@resumewrite@
        {\expandafter\string\csname etaremune@#1\endcsname}%
        {\expandafter\string\csname etaremune@#2\endcsname}}%
      \temp
    \endgroup
  }
  \newcommand*\EM@resumewrite@[2]{% %% Write to aux file
    \immediate\write\@auxout{\xdef#1{\string\noexpand\string\the\numexpr#1+\string\noexpand#2 }}%
  }
\makeatother %% <- change @ back

% A substitute for etaremune:
% \newenvironment{benumerate}[1]{
%     \let\oldItem\item
%     \def\item{\addtocounter{enumi}{-2}\oldItem}
%     \begin{enumerate}
%     \setcounter{enumi}{#1}
%     \addtocounter{enumi}{1}
% }{
%     \end{enumerate}
% }

% \renewcommand*{\cventry}[7][.25em]{
%   \cvitem[#1]{#2}{%
%     {\bfseries#3}%
%     \ifthenelse{\equal{#4}{}}{}{, {\slshape#4}}%
%     \ifthenelse{\equal{#5}{}}{}{, #5}%
%     \ifthenelse{\equal{#6}{}}{}{, #6}%
%     \strut% ROW CHANGED!
%     \ifx&#7&%
%     \else{\newline{}\begin{minipage}[t]{\linewidth}\small#7\end{minipage}}\fi}}



% % bibliography adjustements (only useful if you make citations in your resume, or print a list of publications using BibTeX)
% %   to show numerical labels in the bibliography (default is to show no labels)
% % \makeatletter\renewcommand*{\bibliographyitemlabel}{\@biblabel{\arabic{enumiv}}}\makeatother
% %   to redefine the bibliography heading string ("Publications")
% % \renewcommand{\refname}{Articles}
% % bibliography with mutiple entries
% % \usepackage{multibib}
% % \newcites{book,misc}{{Books},{Others}}
% % Count total number of entries in each refsection
% \AtDataInput{%
%   \csnumgdef{entrycount:\therefsection}{%
%     \csuse{entrycount:\therefsection}+1}}
% 
% % Print the labelnumber as the total number of entries in the
% % current refsection, minus the actual labelnumber, plus one
% \DeclareFieldFormat{labelnumber}{\mkbibdesc{#1}}    
% \newrobustcmd*{\mkbibdesc}[1]{%
%   \number\numexpr\csuse{entrycount:\therefsection}+1-#1\relax}
% 

%% Make years smaller
\let\oldcventry\cventry
\renewcommand{\cventry}[6]{\oldcventry{{\small #1}}{#2}{#3}{#4}{#5}{#6}}

%%
\usepackage{xcolor}
% \AfterPreamble{

% }

%----------------------------------------------------------------------------------
%            content
%----------------------------------------------------------------------------------

\name{Joseph~T.}{Iosue, Ph.D.}
% \title{\normalfont{Doctoral candidate\\[.5em]University of Maryland}}                               % optional, remove / comment the line if not wanted
\title{Senior Researcher}
% \address{Doctoral candidate, University of Maryland}{Joint Center for Quantum Information and Computer Science}{Joint Quantum Institute}% optional, remove / comment the line if not wanted; the "postcode city" and "country" arguments can be omitted or provided empty
\address{Microsoft Quantum}{}{Redmond, Washington}
% \address{3106 Atlantic Bldg.}{University of Maryland}{College Park, MD 20742}% optional, remove / comment the line if not wanted; the
% \phone[mobile]{+1~(234)~567~890}                   % optional, remove / comment the line if not wanted; the optional "type" of the phone can be "mobile" (default), "fixed" or "fax"
% \phone[fixed]{+2~(345)~678~901}
% \phone[fax]{+3~(456)~789~012} 
% \email{jtiosue@umd.edu}
\email{jtiosue@gmail.com}    
\homepage{jtiosue.github.io}                         % optional, remove / comment the line if not wanted
% \social[linkedin]{john.doe}                        % optional, remove / comment the line if not wanted
% \social[twitter]{IosueJoe}                             % optional, remove / comment the line if not wanted
\social[github]{jtiosue}                              % optional, remove / comment the line if not wanted
% \extrainfo{additional information}                 % optional, remove / comment the line if not wanted
% \photo[70pt][0pt]{me.jpg}                       % optional, remove / comment the line if not wanted; '64pt' is the height the picture must be resized to, 0.4pt is the thickness of the frame around it (put it to 0pt for no frame) and 'picture' is the name of the picture file
% \quote{Some quote}                                 % optional, remove / comment the line if not wanted


\begin{document}
\hypersetup{
    % colorlinks=true,
    urlcolor=c1,
    citecolor=blue,
}
%\begin{CJK*}{UTF8}{gbsn}                          % to typeset your resume in Chinese using CJK
%-----       resume       ---------------------------------------------------------
\makecvtitle

\vspace{-5mm}

\section{Education}

% \cventry{A}{B}{C}{D}{E}{F}

\cventry
  {2020 -- 2025}
  {University of Maryland}
  {Doctor of Philosophy (Physics)}
  {JQI Graduate Fellow}
  {}
  {Advisors: Prof.~Alexey Gorshkov and Prof.~Victor Albert}


\cventry
  {2015 -- 2019}
  {Massachusetts Institute of Technology}
  {Bachelor of Science}
  {}
  {}
  {B.S. (Physics), Minor (Computer Science) -- GPA 4.9/5.0}

% \section{Research Highlights}

\cvlistitem{Hello world}

\section{Research Experience}

\cventry
  {2026 --}
  {Microsoft Quantum}
  {Senior Researcher}
  {}
  {}
  {
  Senior researcher on the quantum error correction team at Microsoft
  }

\cventry
  {2025 -- 2026}
  {Johns Hopkins University Applied Physics Laboratory}
  {Senior Scientist}
  {}
  {}
  {
  Quantum research scientist in Research and Exploratory Development (REDD) researching quantum error correction, noise characterization, and sensing
  }


\cventry
  {2020 -- 2025}
  {University of Maryland, Department of Physics}
  {Research Assistant}
  {}
  {}
  {
    Affiliations:
    \begin{itemize}
        \item Joint Quantum Institute (JQI)
        \item Joint Center for Quantum Information and Computer Science (QuICS)
    \end{itemize}
  }


\cventry
  {2024}
  {IBM Quantum Research}
  {Summer intern}
  {}
  {}
  {Supervisor: Kunal Sharma and Minh Tran.}


\cventry
  {2023}
  {Johns Hopkins University Applied Physics Laboratory}
  {Summer intern}
  {}
  {}
  {Supervisor: Dr.~Paraj Titum. Studied a specific quantum sensing and noise characterization problem \cite{iosue2026superresolution}.}


\cventry
  {2019 -- 2020}
  {QC Ware, Corp}
  {Quantum Algorithms Researcher}
  {}
  {}
  {Developed software and algorithms for customer use cases, including \cite{parrish2019a-jacobi-diagon}.}
  

\cventry
  {2018}
  {Los Alamos National Laboratory}
  {Quantum Computing Summer Fellowship}
  {}
  {}
  {Supervisor: Dr.~Patrick Coles. Developed and published a novel quantum algorithm \cite{cirstoiu2020variational-fas}.}


\cventry
  {2017}
  {Joint Quantum Institute, University of Maryland}
  {Summer Researcher}
  {}
  {}
  {Supervisor: Prof.~Alexey Gorshkov. Studied quantum phase transitions via quench dynamics \cite{titum2019probing-ground-}.}


% \cventry
%   {2017 -- 2018}
%   {MIT Laboratory for Nuclear Science}
%   {Undergraduate Researcher}
%   {}
%   {}
%   {Supervisor: Prof.~Or Hen. Studied proton/neutron dynamics in asymmetric nuclei with C\texttt{++} and ROOT.}
  
% \cventry
%   {2016}
%   {MIT Plasma Science and Fusion Center}
%   {Undergraduate Researcher}
%   {}
%   {}
%   {Supervisor: Prof.~Nuno Loureiro. Modeled particle transport in turbulent media using C.}

% \cventry
%   {2015}
%   {MIT Department of Nuclear Science and Engineering}
%   {Undergraduate Researcher}
%   {}
%   {}
%   {Supervisor: Prof.~Emilio Baglietto. Modeled responses of nuclear waste storage canisters to fission waste using finite element software ADINA.}


\cventry
  {2015 -- 2018}
  {Massachusetts Institute of Technology}
  {Undergraduate researcher}
  {}
  {}
  {\begin{itemize}
      \item (2017 -- 2018) Supervisor: Prof.~Or Hen. Studied proton/neutron dynamics with C\texttt{++} and ROOT.
      \item (2016) Supervisor: Prof.~Nuno Loureiro. Modeled particle transport in turbulent media using C.
      \item (2015) Supervisor: Prof.~Emilio Baglietto. Modeled nuclear waste storage canisters and fission waste.
  \end{itemize}}

% \section{Research Interests}


\cvitem{}{Complexity theory and applications to physics, Classical simulation of quantum systems, Verification of quantum computers, Hamiltonian simulation}

\section{Talks}

\cventry
  {}
  {Higher moment theory and learnability of bosonic states}
  {based on \cite{iosue2025higher-moment-t}}
  {}
  {}
  {
  \begin{itemize}
      \item 2026 -- The Theory of Quantum Computation, Communication and Cryptography (TQC), {\em see my slides on \href{https://github.com/jtiosue/learning-bosons/blob/main/TQC2026Presentation/main.pdf}{GitHub}}
      \item 2026 -- Prof.~Scott Aaronson's group meeting
  \end{itemize}
  }

\cventry
  {}
  {Projective toric designs, quantum state designs, and mutually unbiased bases}
  {based on~\cite{iosue2024projective-tori}}
  {}
  {}
  {
  \begin{itemize}
      \item 2025 -- Quantum Information Processing (QIP), {\em see recording on \href{https://youtu.be/zukyDkl9UEY?si=p_ET1ffrKhJnx3vh}{YouTube}}
      \item 2024 -- Codes and Expansions (\href{https://www.math.colostate.edu/~king/codex/}{CodEx}) Seminar (invited, virtual), {\em see recording on \href{https://youtu.be/Wi7-Q98M3-E?si=Rqc28Cn_Ju8ZQmwE}{YouTube}}
  \end{itemize}
  }

\cventry
  {}
  {Continuous-variable quantum state designs:~theory and applications}
  {based on \cite{iosue2024continuous-vari}}
  {}
  {}
  {
    \begin{itemize}
        \item 2023 -- Quantum Information Processing (QIP), {\em see recording on \href{https://youtu.be/yiH0L8PYmTo}{YouTube}}
        \item 2023 -- APS March Meeting
        \item 2022 -- Prof.~David Gross's group seminar (invited, virtual)
        \item 2022 -- CU Boulder journal club (invited, virtual)
        \item 2022 -- University of Maryland JQI-QuICS quantum seminar
        \item 2022 -- APS March Meeting (virtual)
    \end{itemize}
  }


\cventry
  {}
  {Page curves and typical entanglement in linear optics}
  {based on \cite{iosue2023page-curves-and}}
  {}
  {}
  {
    \begin{itemize}
        \item 2023 -- Quantum Algorithms and Applications Collaboratory (\href{https://www.sandia.gov/quantum/quaac/}{QuAAC}) Seminar, Sandia National Laboratory (invited, given by coauthor)
        \item 2023 -- APS March Meeting
    \end{itemize}
  }


% \textit{\small ${}^\dag$: Invited} ~~~ \textit{\small ${}^\ast$: Presented by a co-author}

\section{Publications}

\nocite{*}

\cvitem{}{
    See also:
    \href{https://scholar.google.com/citations?user=JXw0_fEAAAAJ&hl=en}{Google Scholar},
    \href{https://arxiv.org/search/?searchtype=author&query=Iosue+J+T}{arXiv}
}

% \begin{minipage}{13.5mm}
% \end{minipage}
% \hfill
% \begin{minipage}{\textwidth-13.5mm}
%     \printbibliography[heading=none]
% \end{minipage}

\printbibliography[heading=none,env=publist]
% \printbibliography[heading=none]


% \cvitem{}{\small{$^\ddag$ These authors contributed equally.}}
\section{Posters}

\cventry
  {}
  {Page curves and typical entanglement in linear optics}
  {based on \cite{iosue2023page-curves-and}}
  {}
  {}
  {
    \begin{itemize}
        \item 2023 -- Boulder Summer School
        \item 2023 -- Quantum Information Processing (QIP)
    \end{itemize}
  }

\cventry
  {}
  {An initial condition robust outer-loop optimization strategy for QAOA}
  {}
  {}
  {}
  {
    \begin{itemize}
        \item 2019 -- TQC Conference, College Park, Maryland
    \end{itemize}
  }


\section{Selected Projects}


\cventry
  {2019 --}
  {qubovert}
  {Python package (with C extension) for binary optimization}
  {}
  {}
  {
    \begin{itemize}
        \item Created \href{https://github.com/jtiosue/qubovert}{qubovert}, which is particularly designed to aid in converting optimization problems to a form that can be solved with quantum annealers and quantum optimization algorithms.
        \item qubovert can be installed with \texttt{pip install qubovert}, the source code is hosted at \url{github.com/jtiosue/qubovert}, and the documentation is hosted at \url{qubovert.readthedocs.io}.
        \item qubovert currently has \textbf{over 363k downloads} from \href{https://pypi.org/project/qubovert/}{PyPI}, 41 stars and 8 forks on GitHub, and has been \href{https://scholar.google.com/scholar?hl=en&as\_sdt=0\%2C33&q=qubovert&btnG=}{cited} in many research articles.
    \end{itemize}
  }


\cventry
  {2024}
  {(Deep) Q-learning in PyTorch}
  {}
  {}
  {}
  {
  \begin{itemize}
      \item \href{https://github.com/jtiosue/DotsBoxesQLearn}{PyTorch implementation} for self-playing reinforcement learning of a two player game.
  \end{itemize}
  }


\cventry
  {2023}
  {rcal}
  {Python package for review calibration}
  {}
  {}
  {
    \begin{itemize}
      \item Devised a novel review calibration algorithm (\href{https://github.com/jtiosue/rcal/blob/main/report/review_calibration.pdf}{written report})
      \item Implemented the algorithm in a Python \href{https://github.com/jtiosue/rcal}{package}.
    \end{itemize}
  }
    
    
\cventry
  {2019}
  {Powell bounded multivariate optimization}
  {SciPy contribution}
  {}
  {}
  {
    \begin{itemize}
      \item Authored \href{https://github.com/scipy/scipy/pull/10648}{pull request number 10648} on Python’s SciPy package. My contribution is included in the \href{https://github.com/scipy/scipy/releases/tag/v1.5.0}{1.5.0} release and later releases
      \item The pull request implements an additional feature for SciPy’s minimization method. I devised a bounded version of the standard unbounded Powell minimization method and found it to often perform much better than the other gradient-free minimizers. I then implemented this variant in SciPy’s software stack and created the pull request.
    \end{itemize}
  }
    
    
\cventry
  {2018}
  {Quantum Computer Simulator}
  {C\texttt{++} project}
  {}
  {}
  {
    \begin{itemize}
      \item Implemented a \href{https://github.com/jtiosue/Quantum-Computer-Simulator-with-Algorithms}{quantum computer simulator} in C\texttt{++}.
    \end{itemize}
  }



    
\section{Awards and Achievements}


\cventry
  {2023}
  {Boulder Summer School}
  {Non-Equilibrium Quantum Dynamics}
  {University of Colorado, Boulder}
  {}
  {One of 81 accepted into the four-week program}



\cventry
  {2020 -- 2025}
  {JQI Graduate Fellowship}
  {}
  {Joint Quantum Institute}
  {University of Maryland}
  {}


\cventry
  {2019, 2020}
  {NSF GRFP Honorable Mention}
  {}
  {National Science Foundation}
  {}
  {}


% \section{Extra 1}
% \cvlistitem{Item 1}
% \cvlistitem{Item 2}
% \cvlistitem{Item 3. This item is particularly long and therefore normally spans over several lines. Did you notice the indentation when the line wraps?}

% \section{Extra 2}
% \cvlistdoubleitem{Item 1}{Item 4}
% \cvlistdoubleitem{Item 2}{Item 5}
% \cvlistdoubleitem{Item 3}{Item 6. Like item 3 in the single column list before, this item is particularly long to wrap over several lines.}


% \section{Computer skills}
% \cvdoubleitem{category 1}{XXX, YYY, ZZZ}{category 4}{XXX, YYY, ZZZ}
% \cvdoubleitem{category 2}{XXX, YYY, ZZZ}{category 5}{XXX, YYY, ZZZ}
% \cvdoubleitem{category 3}{XXX, YYY, ZZZ}{category 6}{XXX, YYY, ZZZ}

% \input{data/supervision}
\section{Service and Mentoring}

\cventry
  {}
  {Journal and Conference Reviews/Subreviews}
  {}
  {}
  {}
  {
    \begin{itemize}
        \item npj Quantum Information 2026
        \item Quantum journal 2025, 2024, 2023
        \item Quantum Information Processing (QIP) conference 2026, 2025, 2024, 2023
        \item Quantum Computing Theory in Practice (QCTiP) conference 2026
        \item Theory of Quantum Computation, Communication and Cryptography (TQC) conference 2025, 2024
        \item Scientific Reports 2025
    \end{itemize}
  }


\cventry
    {2023 -- 2024}
    {Mentor}
    {}
    {}
    {}
    {Mentored high school summer student Jason Youm studying extensions of my work \cite{iosue2023page-curves-and}. Resulted in publication \cite{youm2025average-renyi-e}. See this \href{https://quics.umd.edu/news/high-school-student-earns-accolades-2023-summer-research-gorshkov-group}{article} about the mentorship.}


\cventry
  {2020 -- 2022}
  {Volunteer tutor}
  {UMD Department of Physics}
  {}
  {}
  {}

\cventry
  {2016}
  {Teaching assistant and grader}
  {MIT Department of Physics}
  {}
  {}
  % {Referred to as ``best TA ever'' by several students in anonymous subject evaluations.}
  {}

% \section{Miscellaneous}


\cventry
  {2020}
  {PHYS.ORG}
  {}
  {Popular science article}
  {\textit{based on \cite{jainQuantumSphericalCodes2023}}}
  {\textit{A framework to construct quantum spherical codes}, \href{https://phys.org/news/2024-06-framework-quantum-spherical-codes.html}{Online}}

\cventry
  {2024}
  {Springer Nature blog}
  {}
  {}
  {\textit{based on \cite{jainQuantumSphericalCodes2023}}}
  {\textit{Polytopes protect Quantum Information}, \href{https://communities.springernature.com/posts/polytopes-protect-quantum-information}{Online}}

\cventry 
  {2024}
  {QuICS blog}
  {}
  {Popular science article}
  {\textit{based on \cite{iosueContinuousvariableQuantumState2022}}}
  {\textit{Carving Up Infinite Quantum Spaces into Simpler Surrogates}, \href{https://quics.umd.edu/news/carving-infinite-quantum-spaces-simpler-surrogates}{Online}}


% \cventry
%   {2023}
%   {On-Line Encyclopedia of Integer Sequences}
%   {based on \cite{iosue2023projective-tori}}
%   {}
%   {}
%   {Our work is included in the \href{https://oeis.org/A108625}{OEIS:A108625} entry}

% \cventry
%   {2022}
%   {On-Line Encyclopedia of Integer Sequences}
%   {based on \cite{iosuePageCurvesTypical2023}}
%   {}
%   {}
%   {Our work is included in the \href{https://oeis.org/A062991}{OEIS:A062991} entry}


\cventry
  {2020}
  {PHYS.ORG}
  {}
  {Popular science article}
  {\textit{based on \cite{cirstoiuVariationalFastForwarding2020}}}
  {\textit{New quantum computing algorithm skips past time limits imposed by decoherence}, \href{https://phys.org/news/2020-10-quantum-algorithm-limits-imposed-decoherence.html}{Online}}

\section{Skills}

\cvitem{}{
  \textbf{Programming}
  \hfill
  Python, C, C\texttt{++}, Mathematica/WolframScript, \LaTeX, Git
}


% \section{Skills \& Qualifications}

% \cvitem{}{
% \begin{itemize}
%     \item PhD in Physics, focus in Quantum Information Science
%     \item 17 publications in the areas of quantum computing, many-body physics, quantum error correction, and machine learning
%     \item Proficient in Python and C, and experienced in C++, including simulation of quantum phenomena
%     \item Experience with developing and applying new problem-solving and numerical methods for quantum phenomena
% \end{itemize}
% }

% \section{References}


\cventry
  {}
  {Alexey Gorshkov}
  {Adjunct Professor}
  {Academic Co-Advisor}
  {\href{mailto:gorshkov@umd.edu}{gorshkov@umd.edu}}
  {National Institute of Standards and Technology (NIST); University of Maryland}


\cventry
  {}
  {Victor Albert}
  {Adjunct Assistant Professor}
  {Academic Co-Advisor}
  {\href{mailto:vva@umd.edu}{vva@umd.edu}}
  {National Institute of Standards and Technology (NIST); University of Maryland}


\cventry
  {}
  {Kunal Sharma}
  {Staff Research Scientist}
  {Internship Advisor and Co-Author}
  {\href{mailto:kunals@ibm.com}{kunals@ibm.com}}
  {IBM T.~J.~Watson Research Center}




\end{document}


